\documentclass[11pt,a4paper]{article}

\usepackage[margin=1in]{geometry}
\usepackage{lmodern}
\usepackage[T1]{fontenc}
\usepackage{amsmath,amssymb}
\usepackage{booktabs}
\usepackage{longtable}
\usepackage{xcolor}
\usepackage{listings}
\usepackage{enumitem}
\usepackage{fancyhdr}
\usepackage[hidelinks,colorlinks=true,linkcolor=blue!55!black,urlcolor=blue!55!black,citecolor=blue!55!black]{hyperref}

% ---- literal underscores etc. in text ----
\usepackage[strings]{underscore}

% ---- breakable DOI links (must follow hyperref) ----
\usepackage{xurl}
\newcommand{\doi}[1]{\href{https://doi.org/#1}{\nolinkurl{#1}}}

% ---- colours ----
\definecolor{codebg}{RGB}{245,245,248}
\definecolor{codekw}{RGB}{0,80,160}
\definecolor{codecm}{RGB}{0,128,64}
\definecolor{codestr}{RGB}{160,60,20}
\definecolor{rule}{RGB}{40,60,120}

% ---- listings ----
\lstdefinestyle{sh}{
  basicstyle=\ttfamily\small,
  backgroundcolor=\color{codebg},
  frame=single, rulecolor=\color{rule!30},
  breaklines=true, columns=fullflexible,
  keepspaces=true, showstringspaces=false,
  keywordstyle=\color{codekw}\bfseries,
  commentstyle=\color{codecm}\itshape,
  stringstyle=\color{codestr},
  xleftmargin=8pt, xrightmargin=4pt, aboveskip=6pt, belowskip=6pt
}
\lstdefinestyle{mat}{style=sh,
  morekeywords={function,end,for,while,if,else,elseif,global,parfor,parfeval,run,addpath,clear}
}
\lstset{style=sh}
\newcommand{\code}[1]{\texttt{#1}}

% ---- headers ----
\pagestyle{fancy}\fancyhf{}
\lhead{\small HEBCPF -- MEX v2.202607}
\rhead{\small User Guide}
\cfoot{\small\thepage}
\renewcommand{\headrulewidth}{0.4pt}

\setlength{\parskip}{4pt}
\setlength{\parindent}{0pt}

% ---- looser line breaking so long code tokens/lists do not overrun the margin ----
\tolerance=3000
\emergencystretch=3em

\begin{document}

% ================= TITLE =================
\begin{titlepage}
\centering
\vspace*{2.2cm}
{\color{rule}\rule{\textwidth}{1.2pt}}\\[8pt]
{\Huge\bfseries HEBCPF}\\[10pt]
{\LARGE Holomorphic-Embedding-Based Continuation}\\[4pt]
{\LARGE Power-Flow All-Solutions Solver}\\[10pt]
{\color{rule}\rule{\textwidth}{1.2pt}}\\[18pt]
{\large\itshape Finding \emph{all} real-valued solutions of the AC power-flow equations}\\[6pt]
{\large\itshape via the ellipsoidal formulation, holomorphic embedding, Pad\'e approximation, and numerical continuation}\\[26pt]

{\large\bfseries Solver line: MEX v2.202607}\\[2pt]
{\large\bfseries Package release: 2026.07.12}\\[4pt]
{\large Windows 64-bit (compiled MEX) build}\\[30pt]

\begin{tabular}{rl}
\textbf{Method:} & Holomorphic Embedding Based Continuation (HEBC)\\[3pt]
\textbf{Originating work:} & D.\ Wu and B.\ Wang, \emph{IEEE Access}, 2019\\[6pt]
\textbf{Foundations:} & Elliptical branch tracing --- B.\ Lesieutre \& D.\ Wu,\\[1pt]
                      & Allerton 2015;\\[3pt]
                      & Algebraic set preserving mappings --- D.\ Wu,\\[1pt]
                      & Ph.D.\ dissertation, UW--Madison, 2017\\[6pt]
\textbf{Document:} & User Guide and Reference Manual\\
\end{tabular}

\vfill
{\small When you publish results obtained with this solver, please cite the HEBC
paper (see Section~\ref{sec:cite}).}\\[6pt]
{\small Package release: July 12, 2026}
\end{titlepage}

\tableofcontents
\newpage

% ================= 1. OVERVIEW =================
\section{Overview}

HEBCPF computes \textbf{all} real-valued solutions of the AC power-flow (load-flow)
equations for a given network, rather than the single operating point returned by a
conventional solver. Knowing the complete solution set is essential for
transient-stability energy-function methods (locating type-1 unstable equilibria),
voltage-stability analysis, and the study of multiple high-voltage operating points.

The solver implements the \textbf{Holomorphic Embedding Based Continuation (HEBC)}
method of Wu and Wang (see Section~\ref{sec:cite}). The power-flow equations are first
recast in an \emph{elliptical form} (Lesieutre and Wu, 2015) so that every
solution-tracing curve is a
bounded loop; the loops are then followed by a hybrid numerical scheme:

\begin{itemize}[leftmargin=1.4em]
  \item \textbf{Holomorphic embedding} traverses the smooth (non-singular) portions of
        a curve in a few very large steps, using power-series / Pad\'e evaluation of the
        holomorphically-embedded voltages;
  \item a classical \textbf{predictor--corrector continuation} takes over only in the
        neighbourhood of singular points (turning points / folds), where it steps
        carefully through the fold and hands control back to the holomorphic routine.
\end{itemize}

Each time the tracing (homotopy) parameter crosses zero, a power-flow solution has been
found. Starting from one known operating point (from MATPOWER's \code{runpf}), the solver
recursively traces the curves emanating from every discovered solution until no new
solution appears.

This release adds a \textbf{consensus Pad\'e-pole estimator} to the classic method
(Section~\ref{sec:pole}), which stabilises the placement of the holomorphic/continuation
hand-off by agreeing a singularity estimate across several buses rather than trusting a
single bus.

Three execution frameworks share the identical numerical hot path:

\begin{center}
\begin{tabular}{@{}lll@{}}
\toprule
Framework & Entry script & Best for\\
\midrule
Serial            & \code{run_batch.m}          & small cases, debugging, reproducibility\\
Parallel (parfor) & \code{run_batch_par.m}      & medium cases, one worker per equation\\
Parallel (parfeval)& \code{run_batch_parfeval.m} & \textbf{large cases (recommended)}\\
\bottomrule
\end{tabular}
\end{center}

% ================= 2. REQUIREMENTS =================
\section{Requirements and dependencies}

\begin{description}[leftmargin=0em,style=nextline]
\item[MATLAB (required).] Developed and tested on Windows 64-bit, MATLAB R2022a and
      later. Uses \code{sparse}, \code{linsolve}, \code{lu}, and standard linear algebra.
\item[MATPOWER (required).] MATPOWER must be installed and on the MATLAB path. The setup
      stage calls MATPOWER functions including \code{runpf} for the initial operating
      point, \code{makeYbus} for the bus-admittance matrix, and
      \code{idx_bus}/\code{idx_brch} for MATPOWER matrix indices. The solver reads the
      bundled MATPOWER-format \code{caseXXX.m} data files, but does \emph{not} ship
      \code{runpf}, \code{makeYbus}, \code{idx_bus}, \code{idx_brch}, or their
      dependencies --- those come from MATPOWER (\url{https://matpower.org}).
\item[Validated setup.] This package release was validated with MATLAB R2022a and
      MATPOWER 7.1. Other MATLAB or MATPOWER versions may work, but should be checked with
      \code{run_merged_case} before production use.
\item[Parallel Computing Toolbox (optional).] Required \emph{only} for the two parallel
      frameworks (\code{run_batch_par.m}, \code{run_batch_parfeval.m}). The serial
      framework runs without it.
\item[Compiled MEX (shipped).] Two hot-path functions are provided as Windows 64-bit
      binaries, \code{Pade_Apprxmt_mex.mexw64} and \code{holomorphic_cont_tri_mex.mexw64}.
      These \code{.mexw64} files are only expected to run on Windows 64-bit MATLAB. They
      make the serial trace roughly $2\times$ faster than the pure-MATLAB build with
      identical results. On Linux/macOS, or after editing hot-path source, you must
      rebuild them (Section~\ref{sec:mex}) with MATLAB Coder and a configured C compiler.
\item[KLU acceleration (shipped).] The bordered Newton solve is accelerated by the
      \code{klurf} MEX (\code{klurf.mexw64}, built from SuiteSparse/KLU), which is shipped and
      enabled automatically; the solver falls back to the dense solve if it is absent
      (Section~\ref{sec:klu}). SuiteSparse/KLU is third-party software under the LGPL-2.1+
      license.
\end{description}

\paragraph{Installing and compiling SuiteSparse/KLU.}
The \code{klurf} MEX links KLU from \textbf{SuiteSparse} (SPDX \code{LGPL-2.1+}). A Windows
64-bit \code{klurf.mexw64} is already bundled; rebuild only for another platform or after
editing the source. You do \emph{not} need a system-wide SuiteSparse install or its CMake
build --- \code{klurf_make.m} compiles just the handful of KLU / AMD / COLAMD / BTF sources
it needs, so a working MATLAB \code{mex} compiler is the only prerequisite.
\begin{enumerate}[leftmargin=1.6em]
  \item Get the SuiteSparse sources (any recent release), e.g.
\begin{lstlisting}[style=sh]
git clone --depth 1 https://github.com/DrTimothyAldenDavis/SuiteSparse
\end{lstlisting}
  \item Copy \code{klurf/klurf_mex.c} and \code{klurf/klurf_make.m} (from the top-level
        \code{klurf/} folder) into \code{SuiteSparse/KLU/MATLAB/}.
  \item In MATLAB, select a C compiler once and run the build from that directory:
\begin{lstlisting}[style=mat]
>> mex -setup C          % pick MSVC / gcc / clang (once per machine)
>> cd '<...>/SuiteSparse/KLU/MATLAB'
>> klurf_make            % compiles klurf.<mexext>, real-valued, CHOLMOD-free
\end{lstlisting}
  \item Copy the resulting \code{klurf.<mexext>} (\code{klurf.mexw64} on Windows,
        \code{klurf.mexa64} on Linux, \code{klurf.mexmaci64}/\code{klurf.mexmaca64} on macOS)
        into this release folder, or anywhere on the MATLAB path. The solver detects it and
        switches on the refactor path automatically.
\end{enumerate}
Further platform notes are in the top-level \code{klurf/README_KLURF.md}.

% ================= 3. INSTALLATION =================
\section{Installation}

\begin{enumerate}[leftmargin=1.6em]
  \item Install MATLAB and MATPOWER; confirm \code{runpf} works
        (\code{which runpf} returns a path).
  \item Unpack this folder (\code{HEBCPF_MEX_v2.20260712}) anywhere on disk.
  \item Add the folder to the MATLAB path each session:
\begin{lstlisting}
>> cd '<path>\HEBCPF_MEX_v2.20260712'
>> addpath(pwd)
\end{lstlisting}
  \item Add only this one HEBCPF release folder to the MATLAB path. Do not add all six
        release folders simultaneously, because they contain many functions and cases with
        the same names.
  \item (Windows x64 only) nothing else is needed --- the shipped \code{.mexw64} binaries
        are picked up automatically. On other platforms, run \code{build_mex}
        (Section~\ref{sec:mex}) or switch to a pure-MATLAB build.
\end{enumerate}

% ================= 4. QUICK START =================
\section{Quick start}

Before running any driver, confirm MATLAB can find MATPOWER:
\begin{lstlisting}[style=mat]
>> which runpf
>> which makeYbus
\end{lstlisting}
Both commands should return paths inside the MATPOWER installation.

From a system shell (non-interactive \code{-batch}):
\begin{lstlisting}[style=sh]
:: SERIAL
matlab -sd "<path>\HEBCPF_MEX_v2.20260712" -batch "addpath(pwd); run('run_batch.m')"

:: PARALLEL (parfor)
matlab -sd "<path>\HEBCPF_MEX_v2.20260712" -batch "addpath(pwd); run('run_batch_par.m')"

:: PARALLEL (parfeval, recommended for big cases)
matlab -sd "<path>\HEBCPF_MEX_v2.20260712" -batch "addpath(pwd); run('run_batch_parfeval.m')"
\end{lstlisting}

From an interactive MATLAB session:
\begin{lstlisting}[style=mat]
>> addpath(pwd)
>> run_batch            % serial batch
>> run_batch_par        % parallel (parfor)
>> run_batch_parfeval   % parallel (parfeval)
>> r = run_merged_case('case14mod','parfeval');  % one case, any mode
>> r.solutions, r.max_residual, r.wall_time
\end{lstlisting}

The two parallel drivers start a local \code{parpool} automatically (default worker
count) if none is running. The one-time pool start-up is excluded from the reported
timing.

% ================= 5. USER OPTIONS =================
\section{User options}

\subsection{Choosing which cases to run}
Each driver has a list near the top:
\begin{lstlisting}[style=mat]
cases_to_run = { 'case9', 'case14mod', 'case39', ... };
\end{lstlisting}
Comment out or add case names. Bundled test systems (expected solution counts, collapsing
antipodal pairs, in parentheses):

\begin{center}\small
\begin{tabular}{@{}llll@{}}
\toprule
\code{case3} (6) & \code{case4gs} (6) & \code{case4BB0} (14) & \code{case4BBc} (12)\\
\code{case5loop} (10) & \code{case5Salam} (10) & \code{case6ww} (6) & \code{case7Salam} (4)\\
\code{case9} (8) & \code{case9Q} (8) & \code{case14mod} (30) & \code{case33bw} (16)\\
\code{case39} (176) & \code{case30}/\code{case_ieee30} (472) & \code{case57mod} (606) & \code{case57} (1322)\\
\bottomrule
\end{tabular}
\end{center}

\subsection{Load scaling}
Each driver contains
\begin{lstlisting}[style=mat]
factor = 1 + 0.19*(0);   % default = 1 (nominal load)
\end{lstlisting}
Changing the multiplier (the ``\code{(0)}'') scales all bus loads and generator outputs,
letting you sweep load level. Leave it at 1 to reproduce the nominal solution counts.

\subsection{Pad\'e pole estimation mode}\label{sec:pole}
The hand-off from the holomorphic routine to the continuation routine is driven by an
estimate of the nearest singularity, taken from the real roots (poles) of the Pad\'e
denominator. This release estimates that pole by \textbf{consensus} across a set of
buses: the true singularity sits at nearly the same magnitude on every bus, whereas
spurious (Froissart) poles are bus-specific. The estimator takes the smallest pole
magnitude agreed on by at least half of the sampled buses (clustered within
\code{pole_cl_tol}), falling back to the global smallest if no consensus forms.

Two options select the bus set (in \code{hybrid_traceloops_4_with_mex.m}, or overridden
serially via the global \code{POLEMODE}):

\begin{description}[leftmargin=0em,style=nextline]
\item[\code{'fixed'} (option 1)] use a predefined bus set \code{param.pole_fixed_buses}
      (shipped default \code{1:2}, auto-filtered to $\le$ bus count).
\item[\code{'random'} (option 2, shipped default)] draw \code{param.pole_nbus} buses
      (default 3) at random \emph{once per curve} and hold them fixed for that whole
      trace. Because this folder seeds the random generator deterministically per
      \code{(solution, equation)} pair (from a solution signature, bus count and equation
      index), the random set is fully reproducible and independent of parallel worker
      scheduling.
\end{description}

To select the mode from a serial session:
\begin{lstlisting}[style=mat]
>> global POLEMODE; POLEMODE = 'fixed';   % or 'random'
>> run_batch
\end{lstlisting}
Parallel workers do not inherit client globals; for the parallel drivers set the default
directly on line~48 of \code{hybrid_traceloops_4_with_mex.m}.

\subsection{Programmatic single-case runs: \code{run_merged_case}}
While the \code{run_batch*} drivers loop over a list of cases and write CSV summaries,
\code{run_merged_case} is the \textbf{unified single-case harness}: it runs \emph{one}
case in \emph{any} of the three frameworks behind a single function call and returns a
structured, machine-readable result instead of printing to the console.
\begin{lstlisting}[style=mat]
[result, solutions] = run_merged_case(case_name, mode, close_pool_on_finish)
%   mode = 'serial' | 'parfor' | 'parfeval'
\end{lstlisting}
It performs the identical preprocessing as the batch drivers (it reuses \code{main.m} as
the setup template, substituting the requested case for the active \code{mpc=case...}
line), traces the case in the selected framework, and \textbf{self-verifies} by
recomputing the maximum residual $|z^\top M_i z - b_i|$ over every solution $z$ and every
equation $i$. The call never throws: on failure it returns a \code{FAIL} result rather
than erroring. The returned \code{result} struct contains:

\begin{center}\small
\begin{tabular}{@{}ll@{}}
\toprule
Field & Meaning\\
\midrule
\code{status} & \code{'PASS'} or \code{'FAIL'}\\
\code{case_name}, \code{mode} & echoed inputs\\
\code{buses} & number of buses\\
\code{solutions} & number of solutions found\\
\code{max_residual} & largest constraint residual over the solution set (correctness check)\\
\code{wall_time} & trace wall-clock time (excludes pool start-up)\\
\code{error_id}, \code{message} & populated only on \code{FAIL}\\
\bottomrule
\end{tabular}
\end{center}
The second output, \code{solutions}, is the full solution matrix \code{Zsave} (each column
a solution). Typical uses: A/B comparison of the three frameworks on one case (identical
setup, compare \code{wall_time}/\code{solutions}/\code{max_residual}); automated
regression or validation checks; and obtaining the solution set and residual directly as
return values. Because it reads \code{main.m}, that file must remain in the folder for
this harness to work.

\subsection{Resuming a long run from a checkpoint}\label{sec:resume}
A large all-solutions run can take hours, so the collection scripts support resuming from a
saved workspace instead of restarting. Each runbook carries a periodic checkpoint hook --- a
commented \code{save} in its ``save temporary data'' block; enable it, or save the full
workspace (the model matrices together with \code{Zsave} and \code{VBook}) explicitly. For
large solution sets, use MATLAB's large-file format:
\begin{lstlisting}[style=mat]
save('mycheckpoint.mat', '-v7.3')
\end{lstlisting}

To resume, \code{load} that file, (re)set \code{wait}, and run the collection script from the
same release folder and with the same framework:
\begin{lstlisting}[style=mat]
load mycheckpoint.mat
wait = 0;
run('runVBook_hybrid_parfeval.m')   % serial: runVBook_hybrid_2023.m
                                    % parfor: runVBook_hybrid_parallel.m
\end{lstlisting}
On start-up the script checks for a \code{VBook} whose row count matches
\code{size(Zsave,2)}; when they agree it sets \code{initbypass} and continues from the last
processed solution rather than re-tracing from the first. Resuming re-seeds the deduplicator
(Section~\ref{sec:dedup}) from the loaded solutions, so a large checkpoint continues at full
speed.

% ================= 6. HOW THE SOLVER WORKS =================
\section{How the solver works}

\subsection{Elliptical formulation}
In rectangular voltage coordinates $x=[V_d^\top\ V_q^\top]^\top$, each power-flow
equation is a quadratic $x^\top C_i x = b_i$. These native surfaces are unbounded
(hyperbolic), so raw branch tracing can run off to infinity. Following
Lesieutre--Wu (2015) and Wu's thesis, a positive-definite \emph{base
ellipsoid} is added to every equation, producing an equivalent system
\[
  x^\top M_i x = 1, \qquad M_i = M_i^\top \succ 0, \quad i=1,\dots,2N_{\text{bus}},
\]
whose zero set is \emph{identical} to the original (algebraic-set preserving) but whose
every equation is now a high-dimensional ellipse centred at the origin. Relaxing one
equation with a tracing parameter $\alpha$ (the code variable \code{aal}) yields a
1-dimensional curve; because the ellipses are compact, each connected component of that
curve is a \textbf{bounded loop} (generically diffeomorphic to a circle), so a trace must
return to its start. Solutions occur wherever $\alpha=0$ along the loop.

\subsection{Holomorphic embedding and Pad\'e}
Rather than crawl the loop with small predictor--corrector steps, HEBC embeds the tracing
parameter into the complex plane and represents each bus voltage as a holomorphic
function, computed as a power series (to degree \code{degree.act}, default 15). All
series degrees share a single LU factorisation, so the coefficients are cheap. A
\textbf{Pad\'e approximant} (balanced numerator/denominator degree, for maximal
convergence radius) then extrapolates the voltages far along the curve in one step, until
the power mismatch exceeds \code{param.mismatch_error}. The real roots of the Pad\'e
denominator reveal the location of the nearest singularity on the curve.

\subsection{Hybrid switching and warm start}
When a holomorphic step approaches a singularity --- signalled either by the corrector
failing to converge, or by the holomorphic increment $|\Delta\alpha|$ collapsing --- the
solver switches to predictor--corrector continuation to pass through the fold. To avoid a
slow ``cold'' warm-up, a \textbf{warm starter} sets the first continuation step to a
fraction of the estimated distance to the singularity (and no larger than a fraction of
the last holomorphic step), then builds a quadratic predictor. After the fold is cleared,
control returns to the holomorphic routine. One holomorphic step costs about
$4\times$ a continuation step but typically covers $\approx 8$--$9$ continuation steps of
arc, which is where the speed-up comes from.

\subsection{Continuation details}
The continuation phase uses a pseudo-arclength predictor (quadratic, from the last three
points), a Newton corrector (\emph{Phase I}), and a variable-rescaling corrector
(\emph{Phase II}) for badly-conditioned near-singular arcs whose secant slope exceeds
\code{param.slope_max}. The step length adapts to the corrector iteration count toward a
target of \code{param.targetsteps}. Each $\alpha$ sign change triggers a solution
extraction and a duplicate check (Section~\ref{sec:dedup}); a trace terminates when it returns to its starting
point (within \code{closure} tolerance) or exhausts the alternation cap
\code{param.outnum}.

\subsection{Linear-solver acceleration: KLU symbolic-once refactor}\label{sec:klu}
The numerical hot path is the \textbf{bordered Newton solve} inside the Phase~I corrector
and inside the solution-resolving routine \code{Resolve_with_mex}. Each is a linear system
whose $N_m\times N_m$ core is the power-flow Jacobian, bordered by one or two dense rows and
columns (the arclength and scaling equations). This solve accounts for roughly
$70$--$98\%$ of the tracing time on large systems. The core, although stored densely by the
classic code, is structurally \emph{sparse}: only about $3\%$ of its entries are nonzero,
reflecting the network connectivity.

The solve is accelerated in three stackable steps, all of which keep the \emph{exact}
Jacobian --- so quadratic Newton convergence and the exact solution set are unchanged:
\begin{enumerate}[leftmargin=1.6em]
  \item \textbf{Sparse core.} \code{MakeJacobianD} returns the Jacobian core directly as a
        sparse matrix, avoiding a dense factorisation and any per-iteration \code{sparse()}
        conversion.
  \item \textbf{Schur complement.} The one or two dense border rows and columns are
        eliminated against the sparse core, leaving only a $1$--$2$ dimensional dense
        system to finish. Concretely, write the bordered Jacobian and right-hand side in
        block form, with the sparse power-flow core $A=J\in\mathbb{R}^{N_m\times N_m}$ in the
        top-left and $m$ thin dense border rows and columns ($m=2$ in the continuation
        corrector --- the arclength and scaling equations --- and $m=1$ in
        \code{Resolve_with_mex}):
        \[
          \begin{bmatrix} A & B\\[2pt] C & D \end{bmatrix}
          \begin{bmatrix} y_1\\[2pt] y_2 \end{bmatrix}
          =
          \begin{bmatrix} f_1\\[2pt] f_2 \end{bmatrix},
          \qquad
          A\in\mathbb{R}^{N_m\times N_m},\;\;
          B\in\mathbb{R}^{N_m\times m},\;\;
          C\in\mathbb{R}^{m\times N_m},\;\;
          D\in\mathbb{R}^{m\times m}.
        \]
        Only $A$ is large and sparse. Eliminating the block $y_1=A^{-1}(f_1-B\,y_2)$ leaves
        the small $m\times m$ \emph{Schur-complement} system for the border unknowns,
        \[
          S\,y_2 = f_2 - C\,g,
          \qquad
          S = D - C\,H,
          \qquad\text{with}\quad
          A\,g = f_1,\;\; A\,H = B,
        \]
        after which $y_1 = g - H\,y_2$ and $dy=-[\,y_1;\,y_2\,]$. The two core solves
        $g=A^{-1}f_1$ and $H=A^{-1}B$ share a \emph{single} factorisation of $A$ and are
        obtained in one call, $\code{GH}=\code{klurf}(A,[\,f_1\ \ B\,])$, with $g$ the first
        column and $H$ the remaining $m$; so the whole bordered solve costs one sparse
        factorisation of $A$ (a cheap \code{klu_refactor} after the first iteration) plus one
        $m\times m$ dense solve, and the dense $(N_m{+}m)\times(N_m{+}m)$ matrix $J_1$ is
        never assembled.
  \item \textbf{KLU symbolic-once refactor.} The sparse core is factorised with \textbf{KLU}
        (SuiteSparse; T.~Davis). Because the sparsity pattern is identical across every
        Newton iteration and every curve of a case, the fill-reducing ordering and symbolic
        analysis are computed \emph{once} (\code{klu_analyze}); each later solve only
        recomputes numeric values on that ordering (\code{klu_refactor}), skipping the
        ordering, symbolic analysis, and pivot search that a fresh factorisation repeats
        every time.
\end{enumerate}

The speed-up grows with the system size $N_m=2N_{\text{bus}}-1$ (measured serially on an
Intel~i9-13900K, MATLAB~R2022a; solution counts identical to the dense solve on every case):
\begin{center}
\begin{tabular}{@{}lccc@{}}
\toprule
System & $N_m$ & full-factor KLU & symbolic-once refactor\\
\midrule
$\le$ case14mod   & $\le 27$ & ${\approx}1.0\times$ & ${\approx}1.0\times$ (overhead-bound)\\
case30 / case33bw & 59 / 65  & ${\approx}1.0\times$ & $1.1$--$1.25\times$\\
case39            & 77       & ${\approx}1.0\times$ & $1.43\times$\\
case118           & 235      & $2.35\times$         & $\mathbf{7.5\times}$\\
\bottomrule
\end{tabular}
\end{center}
Below $N_m\approx 55$ the matrices are small enough that fixed overhead dominates and the
accelerated path is a wash (occasionally a few percent slower); the gain grows monotonically
with $N_m$ and is largest exactly where run time matters. No size gate is applied.

\paragraph{Enabling and overriding.}
The refactor path is provided by the \code{klurf} MEX, built from SuiteSparse/KLU. This
build ships \code{klurf.mexw64} (Windows 64-bit), so acceleration is \textbf{on by default};
if the MEX is absent the solver transparently falls back to the dense solve. To force a mode
for benchmarking, set the global before running:
\begin{center}
\begin{tabular}{@{}ll@{}}
\toprule
\code{global USEKLU; USEKLU = 0;} & dense bordered solve\\
\code{USEKLU = 1;}  & KLU full factorisation each iteration (needs \code{klu.mexw64})\\
\code{USEKLU = 2;}  & KLU symbolic-once refactor (\code{klurf}) --- default when present\\
\code{USEKLU = [];} & auto: refactor if \code{klurf} present, else dense (the default)\\
\bottomrule
\end{tabular}
\end{center}
Correctness is validated by reproducing the exact dense-solve solution counts on every
system up to case39 --- including the ill-conditioned case39 (turning-point slopes
$\sim\!10^9$), where a reused pivot ordering is most at risk.

\subsection{Solution deduplication: keyed bucket lookup}\label{sec:dedup}
As each trace returns candidate solutions, every candidate is checked against the solutions
already found --- different traces routinely re-find the same roots, and only genuinely new
ones are recorded. The classic code did this with a linear scan: each candidate was compared
($L_1$ distance within \code{4e-7}, exact and antipodal) against the \emph{entire} growing
solution set, which is $O(N)$ per candidate and $O(N^2)$ over a run. On a large resumed run
this dominates: at $N\approx2.3\times10^{6}$ solutions the scan costs ${\sim}1$~s per
candidate, so resuming a big checkpoint crawls while a fresh start (small $N$) runs fast.

\code{KeyedDedup.m} replaces the scan with a hash-bucketed lookup. A $1$-Lipschitz key --- the
sum over a few fixed, seeded-random \emph{signature} variables --- buckets the solutions, so a
candidate is compared only against its own bucket and the two neighbouring buckets: $O(1)$
amortised instead of $O(N)$. The match is the same tolerance test on those signature
variables (exact and antipodal). Seeded once from the current solution set, one object serves
the serial, parfor, and parfeval collectors alike. Solution counts are unchanged ---
validated on \code{case9}/\code{case14mod} and on a $2.3\times10^{6}$-solution checkpoint ---
while the per-candidate cost drops from ${\sim}1$~s to ${\sim}40\,\mu$s (${\sim}10^{4}\times$),
which is what makes resuming a large run practical.

% ================= 7. PARAMETERS =================
\section{Parameter reference}

Parameters live in the \code{param} struct built at the top of
\code{hybrid_traceloops_4_with_mex.m} (and mirrored in the parallel drivers). Defaults
below are the shipped values of this legacy-fast + consensus build.

\subsection{Continuation / step control}
\begin{center}\small
\begin{longtable}{@{}p{0.24\textwidth}p{0.14\textwidth}p{0.55\textwidth}@{}}
\toprule
\textbf{Parameter} & \textbf{Default} & \textbf{Meaning}\\
\midrule\endhead
\code{sstep}        & \code{1e-5} & Base arc-length step scale for the continuation.\\
\code{targetsteps}  & \code{3}    & Target number of corrector iterations; drives step-size adaptation.\\
\code{imax}         & \code{40}   & Maximum Newton iterations per corrector step.\\
\code{smax}         & \code{20000}& Maximum continuation steps per \code{branch_trace} call.\\
\code{maximumstep}  & \code{8000} & Maximum step-size multiplier (Phase~I ceiling base).\\
\code{maxphase1}    & \code{maximumstep/2} & Maximum step in Phase~I.\\
\code{maxphase2}    & \code{maxphase1/10}  & Maximum step in Phase~II.\\
\code{minimumstep}  & \code{1e-2} & Minimum step-size multiplier.\\
\code{ef}           & \code{1e-8} & Corrector residual tolerance $\lVert F\rVert$.\\
\code{ex}           & \code{1e-7} & Corrector step tolerance $\lVert \Delta x\rVert$.\\
\code{PS}           & \code{300}  & Phase~II (rescaling) duration, in steps.\\
\code{maxsolu}      & \code{100}  & Maximum solutions collected per traced loop.\\
\code{maxskip}      & \code{50}   & Consecutive numerical stalls tolerated before the loop is skipped.\\
\code{dev}          & \code{20}   & Console print interval (display only).\\
\bottomrule
\end{longtable}
\end{center}

\subsection{Holomorphic / hand-off control}
\begin{center}\small
\begin{longtable}{@{}p{0.24\textwidth}p{0.14\textwidth}p{0.55\textwidth}@{}}
\toprule
\textbf{Parameter} & \textbf{Default} & \textbf{Meaning}\\
\midrule\endhead
\code{mismatch_error}& \code{1e-4} & Maximum holomorphic power-mismatch accepted before shrinking the arc.\\
\code{holonum}      & \code{50}   & Maximum Pad\'e steps per outer iteration before forcing continuation.\\
\code{holo_arc_max} & \code{0.01} & Maximum holomorphic arc length per step.\\
\code{outnum}       & \code{500}  & Cap on holomorphic$\leftrightarrow$continuation alternations per curve (a stall burns this cap).\\
\code{slope_max}    & \code{4e5}  & Maximum secant slope at which a turning point is accepted as passed. Lower (e.g.\ \code{1e4}) is more robust but slower.\\
\code{aal_min}      & \code{1e-6} & Minimum homotopy-parameter increment; the unit for \code{handback}.\\
\code{im_tol}       & \code{1e-3} & Imaginary-part tolerance for accepting a Pad\'e root as a real pole.\\
\code{handback}     & \code{0.1}  & Hand-back hysteresis ($\times$\code{aal_min}): how far past a turning point to travel before returning to the holomorphic routine. \code{0.1} suits case39/case30; \code{0.5} suits case3/case118.\\
\code{interval}     & \code{1}    & Stride for the turning-point finite-difference test.\\
\bottomrule
\end{longtable}
\end{center}

\subsection{Consensus pole estimation}
\begin{center}\small
\begin{longtable}{@{}p{0.24\textwidth}p{0.14\textwidth}p{0.55\textwidth}@{}}
\toprule
\textbf{Parameter} & \textbf{Default} & \textbf{Meaning}\\
\midrule\endhead
\code{pole_mode}       & \code{'random'} & \code{'fixed'} = predefined bus set; \code{'random'} = random set drawn once per curve.\\
\code{pole_fixed_buses}& \code{1:2}      & Bus set used when \code{pole_mode='fixed'} (auto-filtered to $\le$ bus count).\\
\code{pole_nbus}       & \code{3}        & Number of buses drawn when \code{pole_mode='random'}.\\
\code{pole_consensus}  & \code{true}     & If true, take the smallest pole agreed on by $\ge$ half the sampled buses; else global minimum.\\
\code{pole_cl_tol}     & \code{0.10}     & Relative magnitude tolerance for clustering ``the same pole'' across buses.\\
\bottomrule
\end{longtable}
\end{center}

\subsection{Holomorphic degree}
Set in the drivers / \code{main.m}: \code{degree.act = 15} (active power-series degree),
split as \code{degree.numerator = 8}, \code{degree.denominator = 7} for the Pad\'e
approximant. \textbf{The shipped MEX bake these degree values in as compile-time
constants} --- if you change them, regenerate \code{examples.mat} and rebuild the MEX
(Section~\ref{sec:mex}).

% ================= 8. OUTPUTS =================
\section{Outputs}

\textbf{Console:} per-case progress, solution count, and CPU (serial) or wall (parallel)
time, followed by a summary table.

\textbf{CSV summaries} (written into the run folder):
\begin{center}\small
\begin{tabular}{@{}ll@{}}
\toprule
Driver & CSV file\\
\midrule
\code{run_batch.m}          & \code{results_summary.csv}\\
\code{run_batch_par.m}      & \code{results_par_summary.csv}\\
\code{run_batch_parfeval.m} & \code{results_parfeval_summary.csv}\\
\bottomrule
\end{tabular}
\end{center}
Columns: \code{case_name}, \code{cpu_time_sec} (or wall time), \code{n_solutions},
\code{status} (OK / ERROR), \code{error_msg}.
These CSV files are run outputs and are ignored by the repository \code{.gitignore}.

The full solution set is available in the workspace as \code{Zsave} (each column a
solution in the internal variable ordering) after a run; the solver sign-normalises and
deterministically orders it (\code{canonicalize_solutions.m}) so that serial, parfor and
parfeval produce identical output for the same case. Antipodal pairs $\{x,-x\}$ are
identified as one solution.

% ================= 9. FILE MAP =================
\section{File map}

\textbf{Setup / preprocessing:}
\code{makeYbus} from MATPOWER (Ybus), \code{get_quadr_mtrx.m} (power-balance quadratics),
\code{quadr_matrix.m} (per-bus PV/PQ/slack matrices), \code{Solu.m} (initial point from
\code{runpf}), \code{MakeEllipse.m} (base-ellipse generation), \code{Tune_Fac.m}
(continuation scaling), \code{main.m} (reference preprocessing script; its logic is
inlined into the drivers).

\textbf{Core solver (shared hot path):}
\code{hybrid_traceloops_4_with_mex.m} (outer trace-loop orchestrator; builds \code{param},
seeds the RNG, primes the cached operators), \code{holomorphic_hybrid_5_with_mex.m} (one
holomorphic step + consensus pole estimate), \code{branch_trace_hybrid_4_with_mex.m}
(predictor--corrector continuation), \code{Resolve_with_mex.m} (Newton corrector),
\code{holomorphic_cont_tri.m}(+MEX) (Taylor-coefficient triangular solves),
\code{holomorphic_para_sp.m} (sparse $H$ builder), \code{Pade_Apprxmt.m}(+MEX) (Pad\'e
construction), \code{update_V_Pade.m} (Pad\'e evaluation along the arc),
\code{MakeJacobianD.m} (cached sparse Jacobian), \code{quad_form_vals.m} (cached
constraint-residual evaluator), \code{MakeJacobian.m} (scalar Jacobian, used by
\code{Tune_Fac}).

\textbf{Post-processing:} \code{canonicalize_solutions.m} (sign-normalise + deterministic
ordering).

\textbf{Deduplication:} \code{KeyedDedup.m} (bucketed $O(1)$ duplicate lookup shared by all
three collectors; see Section~\ref{sec:dedup}).

\textbf{Parallel housekeeping:} \code{cleanup_parallel_artifacts.m} (clears per-worker
caches, optionally closes the pool to avoid a Windows/R2022a \code{-batch} shutdown
crash).

\textbf{Drivers:}
serial --- \code{runVBook_hybrid_2023.m}, \code{run_batch.m};
parfor --- \code{runVBook_hybrid_parallel.m}, \code{trace_equation_worker.m},
\code{run_batch_par.m};
parfeval --- \code{runVBook_hybrid_parfeval.m}, \code{trace_equation_worker_const.m},
\code{trace_equation_worker.m}, \code{run_batch_parfeval.m};
convenience --- \code{run_merged_case.m} (one case, any mode).

\textbf{Build:} \code{build_mex.m}, \code{examples.mat} (codegen type seeds).

\textbf{Case data:} \code{caseXXX.m} MATPOWER-format test systems.

% ================= 10. MEX =================
\section{Rebuilding the MEX}\label{sec:mex}

The shipped binaries are Windows 64-bit. On other platforms, or after editing any
hot-path source, rebuild with MATLAB Coder:
\begin{lstlisting}[style=mat]
>> cd '<path>\HEBCPF_MEX_v2.20260712'
>> addpath(pwd)
>> build_mex     % rebuilds both .mexw64 from the .m source + examples.mat
\end{lstlisting}
\code{build_mex.m} defines the codegen argument types (\code{degree} as
\code{coder.Constant}; the \code{I.pv}/\code{I.pq} index fields typed variable-size
because \code{case3} has no PV bus). If you change \code{degree.act}, \code{numerator} or
\code{denominator}, regenerate \code{examples.mat} (save the \code{degree}, \code{I} and
\code{Start} structs from any quick run) and rebuild, because the MEX are specialised to
those constant degree values.

\textbf{Windows build gotcha.} If codegen fails with
``\code{Pade_Apprxmt_mex.bat is not recognized}'', the environment variable
\code{NoDefaultCurrentDirectoryInExePath} is set, blocking codegen's generated \code{.bat}.
Unset it \emph{before} launching MATLAB:
\begin{lstlisting}[style=sh]
bash :  unset NoDefaultCurrentDirectoryInExePath ; matlab ...
cmd  :  set "NoDefaultCurrentDirectoryInExePath=" && matlab ...
\end{lstlisting}
Also try \code{mex -setup C} and re-selecting the compiler (developed with MSVC 2019). If
you have no compiler, use the pure-MATLAB build --- same answers, $\sim 2\times$ slower.

To run \emph{without} the MEX (A/B comparison), comment the \code{_mex} call lines in
\code{holomorphic_hybrid_5_with_mex.m} and uncomment the corresponding \code{.m} calls.

% ================= 11. TROUBLESHOOTING =================
\section{Troubleshooting}

\begin{description}[leftmargin=0em,style=nextline]
\item[\code{Undefined function 'runpf'}] MATPOWER is not on the MATLAB path. Install and
      \code{addpath} it.
\item[\code{Undefined function 'makeYbus'}] MATPOWER is not on the MATLAB path, or the
      MATPOWER \code{lib} folder is missing from the path.
\item[\code{Undefined function 'caseXXX'}] Run \code{addpath(pwd)} from this folder; the
      case data files live here.
\item[\code{Undefined function 'Pade_Apprxmt_mex'} / ``generated for a different
      platform''] You are not on Windows x64, or the binary moved. Run \code{build_mex},
      or use a pure-MATLAB build.
\item[\code{codegen ".bat is not recognized"}] See the Windows build gotcha
      (Section~\ref{sec:mex}).
\item[OSQP warning during \code{runpf}] Harmless MATPOWER solver-probe message;
      \code{runpf} still converges (equation error $\sim 0$).
\item[Parallel driver appears to hang at first run] The local \code{parpool} is starting;
      this one-time cost is excluded from the reported timing.
\item[Solution count differs from the expected value] Check the load factor is 1 and that
      the MATPOWER initial solve (\code{runpf}) reported success. Note that the
      \code{pole_mode}/\code{handback} settings affect \emph{tracing efficiency and
      numerical robustness}, not the final solution set on well-behaved cases.
\end{description}

% ================= 12. LICENSE =================
\section{License and third-party notices}

HEBCPF is distributed under the BSD 3-Clause License; see the top-level
\code{LICENSE} file. Some bundled \code{caseXXX.m} files are copied from, derived from,
or modified from MATPOWER and other public test-system sources. Preserve the per-file
attribution notices when redistributing modified versions, and see the top-level
\code{NOTICE} file for third-party attribution details.

The optional KLU acceleration (Section~\ref{sec:klu}) links \textbf{SuiteSparse/KLU}
(Timothy A.~Davis), distributed under the GNU LGPL-2.1-or-later license. The \code{klurf}
MEX source (\code{klurf/klurf_mex.c}) and its build script are included; see
\code{klurf/README_KLURF.md} and the SuiteSparse license terms for details.

% ================= 12. CITATION =================
\section{Citing this solver}\label{sec:cite}

\textbf{If you use results produced by this solver, please cite the HEBC paper:}
\begin{quote}
D.\ Wu and B.\ Wang, ``Holomorphic Embedding Based Continuation Method for Identifying
Multiple Power Flow Solutions,'' \emph{IEEE Access}, vol.~7, pp.~86843--86853, 2019.
doi:\,\doi{10.1109/ACCESS.2019.2925384}.
\end{quote}

BibTeX:
\begin{lstlisting}[style=sh]
@article{wu2019hebc,
  author  = {Wu, Dan and Wang, Bin},
  title   = {Holomorphic Embedding Based Continuation Method for
             Identifying Multiple Power Flow Solutions},
  journal = {IEEE Access},
  volume  = {7},
  pages   = {86843--86853},
  year    = {2019},
  doi     = {10.1109/ACCESS.2019.2925384}
}
\end{lstlisting}

For citing the software package itself, see the top-level \code{CITATION.cff} file.

\subsection*{Further reading}
The method rests on, and is extended by, the following works:

\begin{enumerate}[leftmargin=1.6em]
\item B.\ C.\ Lesieutre and D.\ Wu, ``An Efficient Method to Locate All the Load Flow
      Solutions --- Revisited,'' in \emph{2015 53rd Annual Allerton Conference on
      Communication, Control, and Computing (Allerton)}, Monticello, IL, USA, pp.~381--388,
      IEEE, 2015. doi:\,\doi{10.1109/ALLERTON.2015.7447029}. \emph{(The elliptical
      branch-tracing formulation that underlies this solver.)}
\item D.\ Wu, ``Algebraic Set Preserving Mappings for Electric Power Grid Models and Its
      Applications,'' Ph.D.\ dissertation, University of Wisconsin--Madison, 2017.
      \emph{(The dissertation that formalises the elliptical / algebraic-set-preserving
      formulation and specifies the predictor--corrector continuation machinery --- the
      Phase~I/II correctors, step-length control, and solution/termination logic --- that
      this solver implements.)}
\item D.\ Wu, D.\ K.\ Molzahn, B.\ C.\ Lesieutre, and K.\ Dvijotham, ``A Deterministic
      Method to Identify Multiple Local Extrema for the AC Optimal Power Flow Problem,''
      \emph{IEEE Transactions on Power Systems}, vol.~33, no.~1, pp.~654--668, Jan.\ 2018.
      doi:\,\doi{10.1109/TPWRS.2017.2707925}.
      \emph{(Extends the tracing method to AC optimal power flow.)}
\item B.\ Wang, D.\ Wu, X.\ Su, K.\ Sun, and L.\ Xie, ``A Long-Term Voltage Stability
      Margin Index Based on Multiple Real Power Flow Solutions,'' in \emph{2023 IEEE Power
      \& Energy Society General Meeting (PESGM)}, Orlando, FL, USA, pp.~1--5, IEEE,
      Jul.\ 2023. doi:\,\doi{10.1109/PESGM52003.2023.10252776}.
      \emph{(Defines a voltage-stability margin as the
      minimum voltage-space distance from the high-voltage solution to all other real
      power-flow solutions --- a direct application of this solver's output.)}
\item D.\ Wu and B.\ Wang, ``Influence of Load Models on Equilibria, Stability and
      Algebraic Manifolds of Power System Differential-Algebraic System,'' in \emph{2019
      57th Annual Allerton Conference on Communication, Control, and Computing
      (Allerton)}, Monticello, IL, USA, IEEE, 2019.
      doi:\,\doi{10.1109/ALLERTON.2019.8919649}. \emph{(How load models reshape the
      solution manifold and its equilibria.)}
\item D.\ Wu and B.\ Wang, ``Collection of Numerous Power Flow Solutions of Standard IEEE
      Test Systems,'' \emph{IEEE Dataport}, Jul.\ 2019.
      doi:\,\doi{10.21227/24bh-hj72}. \emph{(The published dataset of reference solution
      sets for the standard IEEE test systems --- useful for validating this solver's
      output.)}
\end{enumerate}

\vfill
\begin{center}\small\itshape
HEBCPF --- MEX v2.202607. This document accompanies the solver release and may be
redistributed with it.
\end{center}

\end{document}
